\documentclass[11pt,a4paper]{article}
\usepackage[margin=27mm]{geometry}
\usepackage{amsmath,amssymb,amsthm,booktabs}
\usepackage{unicode-math}
\setmainfont{STIXGeneral}
\setmathfont{STIX Two Math}
\usepackage[hidelinks]{hyperref}
\newtheorem{theorem}{Theorem}
\title{Regular polytopes with projective special linear groups}
\author{Human Compression project}
\date{}
\begin{document}
\maketitle

The common theory [U] classifies string markings with two proper end
groups, without assuming the intersection property. Here we apply that
theorem to the classification of Leemans and Schulte [LS]. All group and
polytope isomorphisms preserve the ordered distinguished generators.

\begin{theorem}[Leemans--Schulte]
Let $q$ be a prime power. A regular abstract polytope of rank at least
four with full automorphism group $\operatorname{PSL}_2(q)$ has rank four
and belongs to exactly one of the following two classes:
\begin{center}
\begin{tabular}{rlll}
\toprule
$q$ & Polytope & Type & Face vector\\
\midrule
11 & 11-cell & $\{3,5,3\}$ & $(11,55,55,11)$\\
19 & 57-cell & $\{5,3,5\}$ & $(57,171,171,57)$\\
\bottomrule
\end{tabular}
\end{center}
The 11-cell has hemi-icosahedral facets and hemi-dodecahedral vertex
figures; these local maps are interchanged for the 57-cell. Both polytopes
are self-dual and universal for their stated facets and vertex figures.
The edge graph of the 11-cell is $K_{11}$. In both cases the even-word
subgroup is the full automorphism group.
\end{theorem}

\begin{proof}
Let $(r_0,\ldots,r_{n-1})$ be the distinguished string C-marking.
The intersection property makes every subgroup generated by a proper
subset of the marking proper: if $r_i$ belonged to the subgroup
generated by the other generators, their intersection with
$\langle r_i\rangle$ would be both trivial and nontrivial.
In particular both end groups are proper.

The projective restriction and finite classification theorems of [U]
therefore give odd $q$, rank four, and precisely the marked groups
$\operatorname{PSL}_2(11)$, $\operatorname{PSL}_2(19)$ and $S_5$.
The small group $\operatorname{PSL}_2(3)\cong A_4$ is excluded because
its involutions generate only the normal four-group. For the remaining
odd $q$, the group $\operatorname{PSL}_2(q)$ is simple and hence cannot
be $S_5$. Its order $q(q^2-1)/2$ is strictly increasing for positive
$q$, so the other two groups force $q=11$ and $q=19$ respectively.

Conversely, [U] proves the intersection property and ordered-marking
uniqueness for each surviving presentation. The group--polytope
correspondence gives exactly the stated regular polytopes and their
full automorphism groups. The geometry theorem of [U] supplies their
local maps, face vectors, reversal automorphisms, even-word subgroups
and universal covering property. It also proves the edge incidence
statement, rather than inferring it from the face vector alone.
\end{proof}

This recovers [LS, Theorems 1, 3 and 4], the introductory rank
restriction, and the stated geometric conclusions of its Section~5. The complete ten-presentation table of [LS] is proved and
printed once in [U]. The 11-cell and 57-cell retain their attribution
to Gr\"unbaum and Coxeter; the classification above is due to Leemans
and Schulte. Historical subgroup multiplicity formulas and discussion
of other presentations and chiral constructions are outside this
result contract.

The full proofs and finite program of [U] are required dependencies.
Their cost is charged once within this topic. Dickson's classification
and the classical group--polytope and covering correspondences remain
explicit inputs there. The available Lean certificates cover stated
supporting lemmas; they do not formalize this complete classification.

\begin{thebibliography}{9}
\bibitem[U]{U} \emph{Common theory for projective markings and rotations},
companion paper, projective restriction, finite classification and
geometry theorems.
\bibitem[LS]{LS} D. Leemans and E. Schulte,
\emph{Groups of type $L_2(q)$ acting on polytopes},
arXiv:math/0606660v1, the bound fourteen-page author edition;
\emph{Advances in Geometry} \textbf{7} (2007), 529--539.
\end{thebibliography}
\end{document}
